% ---------------------------------------------------------------------
%  1. Introducción
% ---------------------------------------------------------------------
\section{Introducción}
\label{sec:introduccion}

\subsection{Contexto}
El volumen de datos que se almacena y transmite crece de forma sostenida, y con él el costo
de discos, memoria y ancho de banda. La compresión de datos reduce ese costo representando la
misma información con menos bits. Cuando la reconstrucción debe ser exacta ---texto, código
fuente, ejecutables, registros científicos o bases de datos---, solo son aceptables los
métodos \emph{sin pérdida}, en los que la descompresión devuelve bit a bit la entrada
original~\cite{sayood2017,salomon2010}. La teoría de la información establece el límite de lo
alcanzable: ningún código puede representar una fuente con menos bits por símbolo que su
entropía~\cite{shannon1948}.

Dos algoritmos ocupan un lugar central en la historia y en la práctica de la compresión sin
pérdida. La codificación de Huffman~\cite{huffman1952} construye un código prefijo óptimo a
partir de las frecuencias de los símbolos, y el algoritmo Lempel--Ziv--Welch
(LZW)~\cite{welch1984}, derivado de LZ78~\cite{ziv1978}, sustituye frases repetidas por
referencias a un diccionario que se construye durante la lectura. Ambos siguen vigentes: la
codificación de Huffman forma parte de DEFLATE~\cite{rfc1951} y de numerosos formatos, y LZW
se emplea en formatos de imagen y de archivo, y es una referencia habitual en dispositivos con
recursos limitados~\cite{jambek2014,bjaras2019}. Comparar ambos enfoques permite entender dos
formas complementarias de explotar la redundancia: la estadística de los símbolos aislados y la
repetición de secuencias~\cite{bell1989,jayasankar2021}.

\subsection{Problema}
Los estudios comparativos de Huffman y LZW disponibles suelen limitarse a un único dominio de
datos o a archivos de pocos tamaños, reportan tiempos de una sola ejecución y no contrastan
las mediciones con la complejidad teórica de los
algoritmos~\cite{jambek2014,fauzan2022,btoush2008}. En consecuencia, no queda claro en qué
condiciones cada algoritmo comprime mejor, si sus tiempos crecen realmente de forma lineal con
el tamaño de la entrada, ni cómo influyen la entropía de la fuente o el tamaño del diccionario.
El problema que aborda este trabajo es, por tanto: \emph{¿cómo se comparan Huffman y LZW, en
tiempo de ejecución y en tasa de compresión, sobre una gran cantidad de datos de naturaleza
diversa, y en qué medida el comportamiento observado coincide con el análisis teórico?}

\subsection{Objetivos}
\paragraph{Objetivo general.} Comparar teórica y experimentalmente la codificación de Huffman y
el algoritmo LZW en tiempo de compresión, tiempo de descompresión y tasa de compresión,
contrastando las mediciones con el análisis de complejidad de cada algoritmo.

\paragraph{Objetivos específicos.}
\begin{enumerate}[label=\textbf{OE\arabic*.}, leftmargin=*]
  \item Implementar en Java ambos algoritmos (Huffman canónico y LZW de ancho variable de 9 a
        16 bits) y verificar su corrección con una batería de pruebas unitarias que incluya
        ida y vuelta (\emph{round trip}) sin pérdida en el 100\,\% de los casos.
  \item Derivar la complejidad temporal y espacial de codificación y decodificación de cada
        algoritmo y formular hipótesis contrastables a partir de ella.
  \item Medir el tiempo de ejecución en función del tamaño de la entrada desde
        16~KiB ($2^{14}$ bytes) hasta 16~MiB ($2^{24}$ bytes) en cuatro escenarios de datos (hasta 8~MiB en el texto real) y estimar
        el exponente de crecimiento mediante regresión log--log con su coeficiente $R^2$.
  \item Medir la tasa de compresión (bits por byte) frente a la entropía de orden cero en un
        barrido controlado de alfabetos y en los corpus estándar Canterbury y Silesia.
  \item Cuantificar el efecto del tamaño máximo del diccionario de LZW (de 9 a 16 bits) sobre la
        tasa de compresión y el tiempo.
\end{enumerate}

\subsection{Preguntas de investigación}
\begin{enumerate}[label=\textbf{RQ\arabic*.}, leftmargin=*]
  \item ¿El tiempo de codificación y decodificación de Huffman y de LZW crece con el tamaño de
        la entrada según lo que predice su análisis de complejidad?
  \item ¿Cómo se relaciona la tasa de compresión de cada algoritmo con la entropía de orden cero
        y con el tipo de dato (fuentes sin memoria frente a datos con frases repetidas)?
  \item ¿Qué efecto tiene el tamaño máximo del diccionario de LZW sobre la tasa de compresión y
        el tiempo de ejecución?
\end{enumerate}

\subsection{Contribuciones}
\begin{itemize}
  \item Dos implementaciones documentadas y probadas con una interfaz común
        (\texttt{Codec}): Huffman canónico estático y LZW de ancho variable con reinicio del
        diccionario.
  \item Un análisis de complejidad de ambos algoritmos ligado a las decisiones concretas de
        implementación (montículo binario, decodificación canónica, tabla hash de
        direccionamiento abierto).
  \item Un protocolo experimental reproducible, con calentamiento del JIT, repeticiones,
        medianas e intervalos de confianza~\cite{georges2007}, y ajuste de pendientes log--log
        para validar la complejidad~\cite{mcgeoch2012}.
  \item Una evaluación sobre datos sintéticos controlados y sobre los corpus de referencia
        Canterbury~\cite{arnold1997} y Silesia~\cite{deorowicz2003}, con DEFLATE como punto de
        comparación externo.
\end{itemize}

\subsection{Organización del documento}
La \cref{sec:estado-arte} revisa críticamente el estado del arte: los fundamentos teóricos del
problema, las familias de algoritmos que lo resuelven y los estudios comparativos previos. La
\cref{sec:marco-teorico} presenta el marco teórico: definiciones, ambos algoritmos con su
pseudocódigo, ejemplos resueltos, el análisis de complejidad y las hipótesis H1--H4. La
\cref{sec:metodologia} describe gráficamente la metodología, la implementación, los conjuntos
de datos y el protocolo de medición. La \cref{sec:resultados} reporta los resultados de los
experimentos E1--E4, y la \cref{sec:discusion} los contrasta con la teoría. La
\cref{sec:conclusiones} responde las preguntas de investigación y plantea el trabajo futuro.
Los anexos detallan cómo reproducir el estudio y el detalle de la implementación de cada componente.
